\section{Architettura del Codice}\label{architettura-del-codice}

A livello software, BlueAgent organizza l'interazione con l'host
disaccoppiando l'API in \passthrough{\lstinline!/api!} dalla gestione a
basso livello dei container e delle componenti di rete, gestita ad alto
livello dai servizi in \passthrough{\lstinline!/services!}, che contiene
la logica di business, mentre le complessità operative sono delegate a
moduli e driver specializzati (come \passthrough{\lstinline!/engine!},
\passthrough{\lstinline!/network!}, \passthrough{\lstinline!/verifier!},
\passthrough{\lstinline!/images!} e
\passthrough{\lstinline!/lifecycle!}).

\subsection{\texorpdfstring{Astrazione del Runtime:
\texttt{ContainerEngine}}{Astrazione del Runtime: ContainerEngine}}\label{astrazione-del-runtime-containerengine}

BlueAgent non invoca direttamente il demone Incus all'interno della
logica delle challenge. Tutte le operazioni di gestione del ciclo di
vita dei container sono intermediate da un'interfaccia astratta
denominata \passthrough{\lstinline!ContainerEngine!}.

Questa classe astratta formalizza esclusivamente le primitive necessarie
ai requisiti operativi del Cyber Range:

\begin{itemize}
\tightlist
\item
  Creazione, avvio, arresto forzato ed eliminazione delle istanze.
\item
  Duplicazione rapida e gestione delle interfacce di rete.
\item
  Esecuzione dei comandi nei container.
\end{itemize}

Disaccoppiare l'interfaccia dall'implementazione permette ai servizi di
non dover conoscere come avviene la gestione dei container, e di
predisporre il codice per poter utilizzare altri gestori.

L'implementazione concreta \passthrough{\lstinline!IncusEngine!}
comunica con il demone di sistema tramite Unix Domain Socket
(\passthrough{\lstinline!/var/lib/incus/unix.socket!}), traducendo le
primitive astratte in chiamate REST HTTP asincrone.

\begin{quote}
\passthrough{\lstinline!Unix Domain Socket (UDS)!}: Meccanismo di
Inter-Process Communication (IPC) gestito interamente dal kernel del
sistema operativo. La comunicazione avviene trasferendo i dati
direttamente tra i buffer di memoria dei processi nello spazio kernel
senza attraversare lo stack di rete. L'endpoint è rappresentato da un
file speciale sul filesystem (es:
\passthrough{\lstinline!/var/lib/incus/unix.socket!}).\footnote{\label{fn:ref_21}Linux
  manual page, unix(7) - Sockets for local interprocess communication:
  https://man7.org/linux/man-pages/man7/unix.7.html}
\end{quote}

\subsection{\texorpdfstring{Ciclo di Vita della Sandbox: Async Context
Manager
(\texttt{OVNVerifier})}{Ciclo di Vita della Sandbox: Async Context Manager (OVNVerifier)}}\label{ciclo-di-vita-della-sandbox-async-context-manager-ovnverifier}

Il processo di verifica richiede di istanziare una replica fedele della
rete di produzione che deve esistere esclusivamente per la durata dei
test. In caso di qualsiasi tipo di fallimento, l'infrastruttura clonata
deve essere rilasciata per evitare risorse allocate e sprecate.

Per garantire questa proprietà, il ciclo di vita della sandbox è
incapsulato all'interno della classe
\passthrough{\lstinline!OVNVerifier!}, implementata come
\textbf{Asynchronous Context Manager} di Python.

\begin{quote}
\textbf{Pattern RAII (Resource Acquisition Is Initialization):}
L'allocazione delle risorse è vincolata dall'acquisizione del contesto,
mentre il loro rilascio e smantellamento deterministico è garantito
all'uscita dal blocco, gestendo automaticamente qualsiasi anomalia o
interruzione. \footnote{\label{fn:ref_22}Bjarne Stroustrup, Resource Acquisition Is
  Initialization (RAII), C++ Programming Language.}
\end{quote}

Il ciclo di vita dell' \passthrough{\lstinline!OVNVerifier!} segue tre
fasi:

\begin{enumerate}
\def\labelenumi{\arabic{enumi}.}
\tightlist
\item
  \textbf{Inizializzazione (\passthrough{\lstinline!\_\_aenter\_\_!}):}

  \begin{itemize}
  \tightlist
  \item
    Creazione della rete OVN temporanea
  \item
    Duplicazione dei container tramite snapshot ZFS
  \item
    Riconnessione dell'interfaccia \passthrough{\lstinline!eth0!} dei
    cloni alla nuova rete OVN
  \item
    Istanziazione opzionale del container \emph{Verifier Runner}.
  \end{itemize}
\item
  \textbf{Esecuzione delle Verifiche
  (\passthrough{\lstinline!run\_checks!}):}

  \begin{itemize}
  \tightlist
  \item
    Esecuzione controllata dei comandi di test con timeout
  \end{itemize}
\item
  \textbf{Rilascio e Teardown
  (\passthrough{\lstinline!\_\_aexit\_\_!}):}

  \begin{enumerate}
  \def\labelenumii{\arabic{enumii}.}
  \tightlist
  \item
    Arresto ed eliminazione del container
    \passthrough{\lstinline!verify-runner!};
  \item
    Arresto forzato ed eliminazione di tutti i container cloni;
  \item
    Eliminazione della rete virtuale OVN.
  \end{enumerate}
\end{enumerate}

L'uso dell'\textbf{Async Context Manager} fornisce dunque un metodo
sicuro per utilizzare l'oggetto \passthrough{\lstinline!OVNVerifier!}.

\subsection{Concorrenza e Two-Tier
Locking}\label{concorrenza-e-two-tier-locking}

Incus non supporta operazioni concorrenti sul medesimo container.
Inoltre per evitare problemi di accesso alle risorse è opportuno
implementare un sistema di lock.

BlueAgent implementa un meccanismo di sincronizzazione a due livelli
(\passthrough{\lstinline!TwoTierLock!}), dove sequenzialmente si
richiede un tier dopo l'altro:

\begin{itemize}
\tightlist
\item
  \textbf{Tier 1 (Intra-processo):} utilizza
  \passthrough{\lstinline!asyncio.Lock!} per effettuare il lock delle
  risorse all'interno di un singolo processo.
\item
  \textbf{Tier 2 (Inter-processo):} applica un lock non bloccante sul
  filesystem (\passthrough{\lstinline!fcntl.flock!}) condiviso tra più
  processi
\end{itemize}

Questo lock protegge l'accesso alle risorse in due casi:

\begin{itemize}
\tightlist
\item
  BlueAgent esegue su un solo processo (worker), ma le richieste
  potrebbero essere gestite da più task contemporaneamente. Allora il
  lock Tier 1 protegge l'accesso alla stessa risorsa da parte di due
  task sullo stesso worker.
\item
  BlueAgent esegue su più processi (worker), allora il lock Tier 2
  protegge l'accesso alla stessa risorsa da parte di due task su worker
  differenti.
\end{itemize}

\begin{quote}
L'esecuzione multi-processo è abilitata configurando il numero di worker
e viene gestito da uvicorn.
\end{quote}

\begin{quote}
\passthrough{\lstinline!uvicorn!}: è il web server ASGI ad alte
prestazioni su cui si appoggia FastAPI; mentre FastAPI modella e valida
le rotte dell'applicazione, Uvicorn governa l'accettazione delle
connessioni HTTP e ne distribuisce il carico tra worker e task
asincroni.
\end{quote}

La risorsa principale su cui avviene la mutua esclusione nel codice è il
team.

Il lock può essere effettuato anche sulla singola istanza, o a livello
globale, utile quando si interagisce con meccanismi di sistema come
\passthrough{\lstinline!nftables!}.

I lock inoltre sono rientranti, dunque uno stesso task può richiedere
più volte lo stesso lock incrementando all'acquisizione il suo
contatore.

\subsection{Pipeline Asincrona non Bloccante su Socket
Unix}\label{pipeline-asincrona-non-bloccante-su-socket-unix}

BlueAgent utilizza il modello asincrono per comunicare con la socket
Incus ed adattarsi al metodo \passthrough{\lstinline!wait!} che espone.
La classe \passthrough{\lstinline!IncusEngine!} comunica con il socket
UDS locale (\passthrough{\lstinline!/var/lib/incus/unix.socket!}) per
mezzo di \passthrough{\lstinline!httpx.AsyncClient!}.

Successivamente se Incus gestisce la richiesta in background, BlueAgent
interroga l'endpoint di Incus
\passthrough{\lstinline!/1.0/operations/<uuid>/wait!} per sapere quando
l'operazione viene rilasciata in modo asincrono.

Inoltre, l'esecuzione dei comandi sulla macchina host avviene tramite
\passthrough{\lstinline!asyncio.create\_subprocess\_exec!} all'interno
del runner asincrono, evitando chiamate bloccanti sul task.

Evitare il busy waiting (o comunque operazioni bloccanti) sui task
uvicorn è essenziale in BlueAgent per mantenere il sistema reattivo
nella gestione delle richieste che spesso possono richiedere vari
secondi per essere soddisfatte.

\subsection{\texorpdfstring{Resilienza e Riconciliazione dello Stato
(\texttt{LifecycleManager})}{Resilienza e Riconciliazione dello Stato (LifecycleManager)}}\label{resilienza-e-riconciliazione-dello-stato-lifecyclemanager}

Qualora l'host subisca un riavvio forzato improvviso o il processo di
BlueAgent venga terminato bruscamente con
\passthrough{\lstinline!SIGKILL!}, le routine di pulizia
(\passthrough{\lstinline!\_\_aexit\_\_!}) non possono intervenire,
lasciando container e switch virtuali orfani sul sistema.

Per risolvere queste anomalie BlueAgent ha un demone in background, il
\passthrough{\lstinline!LifecycleManager!}. Ogni 600 secondi (10 minuti,
o comunque un tempo configurabile) il demone interroga i metadati di
Incus, per ogni container di verifica, ed elimina quelli con età
superiore al \textbf{grace period} (per esempio, 600 secondi).

In questo modo si eliminano le risorse di verifica rimaste orfane per
qualsiasi motivo.

\subsection{Compilazione delle
immagini}\label{compilazione-delle-immagini}

La generazione delle immagini prevede di calcolare lo \textbf{SHA-256}
della configurazione ricevuta. Questo permette di gestire i seguenti
casi:

\begin{itemize}
\tightlist
\item
  Se un'immagine con stesso hash e stesso id è già presente nel database
  allora il sistema non fa nulla.
\item
  Se l'hash differisce ma l'id è presente nel database allora si può
  forzare la ricompilazione dell'immagine usando la nuova
  configurazione.
\end{itemize}

La compilazione delle immagini è protetta da un lock sul nome
dell'immagine, ed il log relativo viene riversato su un file
\passthrough{\lstinline!.log!} accessibile tramite
\passthrough{\lstinline!GET /images/\{alias\}/logs!}.
